%% 4-9 波をたし合わせる：加法定理とうなり（第13回：組 V4-5。ページ制限なし。偶数ページ）
%% ほかの巻が指している形：加法定理 cos(a+b)=cos a cos b − sin a sin b、sin(a+b)=sin a cos b + cos a sin b（5-15）、
%%   2倍角 cos2θ=1−2sin²θ=2cos²θ−1、sin2θ=2sinθcosθ、sin²θ=(1−cos2θ)/2、cos²θ=(1+cos2θ)/2（5-15・第5巻の巻末）、
%%   合成 a sinθ + b cosθ = √(a²+b²) sin(θ+α)（合成の角は α。φ は第5巻で初期位相に使う）、sin2θ で飛距離が45°で最大（物理シリーズ）
\makeatletter
% 小見出し（太字の1行）
\newcommand\bkivjhead[1]{\par\addvspace{6pt}\Needspace{6\baselineskip}\noindent{\sffamily\bfseries\color{\bkTT}#1}\par\nobreak\vspace{2pt}}
% 囲みの中の小見出し（それだけで1行）
\newcommand\bkivjbox[1]{\par\addvspace{3pt}\noindent\textbf{#1}\par\noindent\ignorespaces}
% 大事な式のわく
\newcommand\bkivjframe[2][100mm]{\par\addvspace{4pt}{\centering\tikz\node[draw=\bkTT, line width=0.9pt, fill=bkTheme!10, rounded corners=1.5mm, inner xsep=6mm, inner ysep=2.5mm, align=center]{#2};\par}\addvspace{4pt}\noindent\ignorespaces}
% 例題の途中式の行（折り返した行も、そろえた位置から始まる）
\newcommand\bkivjq{\par\vspace{2pt}\noindent\hangindent=4.3\zw\hangafter=1\hspace*{2.6\zw}}
\newcommand\bkivja{\par\vspace{2pt}\noindent\hangindent=2.6\zw\hangafter=0\relax}
\newcommand\bkivjb{\par\vspace{2pt}\noindent\hangindent=4.3\zw\hangafter=0\relax}
% 座標の軸（原点中心）
\newcommand\bkivjaxes{%
  \draw[bkGray, line width=0.6pt, -{Stealth[length=4pt]}] (-25,0) -- (25,0) node[anchor=west, font=\sffamily\footnotesize, text=bkInk] {$x$};
  \draw[bkGray, line width=0.6pt, -{Stealth[length=4pt]}] (0,-24) -- (0,25) node[anchor=south, font=\sffamily\footnotesize, text=bkInk] {$y$};}
% 式を、まん中に1行で置く
\newcommand\bkivjline[1]{\par\addvspace{3pt}{\centering #1\par}\addvspace{3pt}\noindent\ignorespaces}
\makeatother

\bkchapter{4}{9}{波をたし合わせる：\\加法定理とうなり}{\bklevel{4}}{波をたし合わせる：加法定理とうなり}
\bkmichishirube{3-9,4-4,4-7,4-8}{5-15,6-10,9-1}

\begin{goal}
\begin{itemize}[label={\color{\bkTT}●}, leftmargin=1.6\zw, itemsep=0pt, topsep=1pt]
\item \bklevel{4}\enspace 加法定理を単位円で導き、$\sin(a\pm b)$・$\cos(a\pm b)$ を計算できる。
\item \bklevel{4}\enspace 2倍角の式と $\sin^2\theta=\dfrac{1-\cos2\theta}{2}$ を作り、使える。
\item \bklevel{4}\enspace 合成と「積と和の式」で、波のたし算（うなり・電力）を説明できる。
\end{itemize}
\end{goal}

\begin{hatena}
\noindent\begin{minipage}[c]{0.55\linewidth}
\raggedright
ギターの2本の弦を、ほとんど同じ高さの音で同時に鳴らすと、音が「ウワン、ウワン」と、大きくなったり小さくなったりします。これを\textbf{うなり}といいます。2つの音の高さが近づくほど、ウワンウワンはゆっくりになり、ぴったり同じになると消えます。

2つの音を重ねると、なぜ音の大きさが、くり返し変わるのでしょう。1秒に何回ウワンと聞こえるかは、何で決まるのでしょう。
\end{minipage}\hfill
\begin{minipage}[c]{0.42\linewidth}
\centering
\begin{tikzpicture}
\begin{axis}[width=58mm, height=36mm, axis lines=none, xmin=0, xmax=1.02, ymin=-2.3, ymax=2.6, clip=false]
  \addplot[bkBlue, line width=0.9pt, domain=0:1, samples=500] {sin(360*10*x)+sin(360*12*x)};
  \addplot[bkOrange, line width=0.9pt, dash pattern=on 2pt off 1.5pt, domain=0:1, samples=100] {2*abs(cos(360*x))};
  \addplot[bkOrange, line width=0.9pt, dash pattern=on 2pt off 1.5pt, domain=0:1, samples=100] {-2*abs(cos(360*x))};
  \draw[bkGray, line width=0.6pt] (axis cs:0,0) -- (axis cs:1.02,0);
  \foreach \x/\l in {0/大,0.25/小,0.5/大,0.75/小,1/大} {\edef\tmp{\noexpand\node[font=\noexpand\sffamily\noexpand\footnotesize, anchor=south, inner sep=1pt] at (axis cs:\x,2.15) {\l};}\tmp}
\end{axis}
\end{tikzpicture}
\end{minipage}
\end{hatena}

\noindent 4-8では、1つの回る点の影、つまり1つの波を調べました。音・光・電気の世界では、いくつもの波が重なり合います。この章では、2つの波をたし合わせると何が起きるかを調べます。そのための道具が「加法定理」です。

%% ---------------- やってみよう・思い出そう ----------------
\begin{yattemiyou}[表計算で、ずれた2つの波をたそう]
\begin{enumerate}[label={\bkcircled[bkGreen]{\arabic*}}, leftmargin=2.4\zw, itemsep=1pt, topsep=0pt]
\item A列に、角 $\theta$ を0から360まで、30ずつ入れよう。
\item B列に \texttt{=SIN(RADIANS(A2))} と入れて、$\sin\theta$ を計算しよう（4-2：表計算の SIN はラジアンで計算する）。
\item C列に \texttt{=SIN(RADIANS(A2+60))} と入れよう。60°先に進んだ波 $\sin(\theta+60°)$ です。
\item D列に \texttt{=B2+C2} と入れて、2つの波をたそう。B・C・D列をグラフにしよう。たした波は、どんな形かな？　山の高さは $1+1=2$ になるかな？
\end{enumerate}
\begin{bkcenter}
\bksheet{4}{24}{%
1 & $\theta$ & $\sin\theta$ & $\sin(\theta+60°)$ & たした値 \\
2 & 0 & 0 & 0.866 & 0.866 \\
3 & 30 & 0.5 & 1 & 1.5 \\
4 & 60 & 0.866 & 0.866 & 1.732 \\
5 & 90 & 1 & 0.5 & 1.5 \\
6 & 120 & 0.866 & 0 & 0.866 \\
7 & 150 & 0.5 & $-0.5$ & 0 \\
8 & 180 & 0 & $-0.866$ & $-0.866$ \\}
\bktabcaption{180°から先は、180°までの値にマイナスをつけたものになる（4-8 の $\sin(\theta+180°)=-\sin\theta$）。}
\end{bkcenter}
\end{yattemiyou}

\begin{bkfigure}
\begin{tikzpicture}
\begin{axis}[bkplot, width=118mm, height=50mm, xmin=0, xmax=395, ymin=-2.05, ymax=2.25, xtick={90,180,270,360}, xticklabels={}, ytick={-1.732,-1,1,1.732}, yticklabels={$-1.73$,$-1$,1,1.73}, xlabel={$\theta$}, tick label style={font=\sffamily\footnotesize}]
  \addplot[bkcurve, line width=1.1pt, domain=0:360, samples=120] {sin(x)};
  \addplot[bkcurve2, line width=1.1pt, domain=0:360, samples=120] {sin(x+60)};
  \addplot[bkcurve3, domain=0:360, samples=120] {sin(x)+sin(x+60)};
  \node[font=\sffamily\footnotesize, anchor=north, inner sep=1pt] at (axis cs:90,-2.07) {90°};
  \node[font=\sffamily\footnotesize, anchor=north, inner sep=1pt] at (axis cs:180,-2.07) {180°};
  \node[font=\sffamily\footnotesize, anchor=north, inner sep=1pt] at (axis cs:270,-2.07) {270°};
  \node[font=\sffamily\footnotesize, anchor=north, inner sep=1pt] at (axis cs:360,-2.07) {360°};
  \node[font=\sffamily\footnotesize, text=bkGreen!70!black, anchor=south west, inner sep=1pt] at (axis cs:66,1.78) {たした波（山の高さ1.73）};
\end{axis}
\end{tikzpicture}
\bkfigcaption{青が $\sin\theta$、オレンジが $\sin(\theta+60°)$、緑が2つをたした波。}
\end{bkfigure}

\noindent たした波（緑）も、同じ周期360°の、なめらかな波になりました。でも、山の高さは2ではなく、約1.73です。2つの波の山の位置が60°ずれているので、山と山がぴったりは重ならないからです。山の位置は60°で、もとの $\sin\theta$ の山（90°）より30°早くなっています。この1.73と30°は、どこから出てくるのでしょう。この章の⑥で答えます。

\begin{omoidasou}
\bkrecall{4-4}{2点 $(x_1,\ y_1)$、$(x_2,\ y_2)$ の距離 $d$ は、$d^2=(x_2-x_1)^2+(y_2-y_1)^2$。}
\bkrecall{4-7}{余弦定理：三角形で、角 $C$ をはさむ2辺が $a$、$b$ なら、向かいの辺 $c$ は\\$c^2=a^2+b^2-2ab\cos C$。}
\bkrecall{4-8}{単位円の上の点は $(\cos\theta,\ \sin\theta)$。$\sin^2\theta+\cos^2\theta=1$。\\$\cos(-\theta)=\cos\theta$、$\sin(-\theta)=-\sin\theta$。\\90°回すと $(x,\ y)\to(-y,\ x)$。$\cos(\theta+90°)=-\sin\theta$、$\sin(\theta+90°)=\cos\theta$。}
\bkrecall{3-8}{展開：$(p-q)^2=p^2-2pq+q^2$。}
\bkrecall{3-9}{投げ上げたボールの高さ：\qty{20}{m/s}で真上に投げると、$t$ 秒後に $20t-5t^2$（m）。}
\bkrecall{3-10}{家のコンセントの「\qty{100}{V}」は実効値。電圧の山の高さは $100\sqrt2\approx\qty{141}{V}$。}
\end{omoidasou}

\noindent この章でも、「≈」（ほぼ等しい。5-8 でくわしく学びます）を使います。

%% ---------------- 数学の道具 ----------------
\Needspace{14\baselineskip}
\bksection{doguu}{数学の道具}

\bksubsection{1}{$\sin(a+b)$ は、$\sin a+\sin b$ ではない}

\leavevmode\label{bk:4-9:hajime}やってみようでは、$\sin\theta$ と $\sin(\theta+60°)$ をたしました。この和を計算するには、$\sin(\theta+60°)$ のように「2つの角の和の $\sin$」を、$\theta$ の $\sin$ と $\cos$ で表す式がほしくなります。

まず、思いつきやすい考えを、数で確かめます。
\bkng{$\sin(a+b)=\sin a+\sin b$}
\noindent $a=30°$、$b=60°$ とすると、左は $\sin90°=1$、右は $\sin30°+\sin60°$ で、$0.5+0.866=1.366$。ちがう数です。そもそも $\sin$ の値は1をこえないので、右のように1.366になることはありません。
\par\vspace{6pt}
\noindent\begin{minipage}[c]{0.55\linewidth}
\raggedright
図で考えると、わけがわかります。$\sin$ は単位円の上の点の高さです。角を30°から90°へと3倍にしても、高さは0.5から1へと2倍にしかなりません。円の上の点は、上へ行くほど横へ回りこむので、角と高さは比例しないのです（4-8⑥：上の近くでは、影はゆっくり動く）。

だから、$\sin(a+b)$ を正しく計算するきまりを、単位円で作ります。それが\textbf{加法定理}（かほうていり）です。「加法」はたし算のことです。
\end{minipage}\hfill
\begin{minipage}[c]{0.4\linewidth}
\centering
\begin{tikzpicture}[x=1mm, y=1mm]
  \draw[bkGray, line width=0.6pt] (-6,0) -- (26,0);
  \draw[bkGray, line width=0.6pt] (0,-3) -- (0,27);
  \draw[bkInk!45, line width=0.8pt] (-3:22) arc (-3:100:22);
  \foreach \a/\c in {30/bkBlue,90/bkOrange} {\draw[bkInk, line width=0.8pt] (0,0) -- (\a:22); \fill[\c] (\a:22) circle (1.2);}
  \draw[bkBlue, line width=1.6pt] ({22*cos(30)},0) -- (30:22);
  \node[font=\sffamily\footnotesize, anchor=west, text=bkBlue!80!black] at ({22*cos(30)+0.8},5.5) {0.5};
  \node[font=\sffamily\footnotesize, anchor=west, text=bkOrange!80!black] at (1.6,22) {1};
  \node[font=\sffamily\footnotesize] at (15:8) {30°};
  \node[font=\sffamily\footnotesize] at (62:9) {90°};
\end{tikzpicture}
\end{minipage}

\bksubsection{2}{加法定理を作る：道すじごと回す}

\leavevmode\label{bk:4-9:kahou}単位円の上で、角 $b$ の所にある点Qは、原点Oから「横に $\cos b$ 進み、そこから上に $\sin b$ 進む」と着く点です（4-8②）。この道すじを、\textbf{まるごと角 $a$ だけ回します}。回したあとの点Rは、角 $a+b$ の所にあるので、Rの座標は $(\cos(a+b),\ \sin(a+b))$ です。いっぽう、回した道すじを1本ずつ調べても、Rの座標が出せます。2つを等しいとおけば、式ができます。

\begin{bkfigure}
\begin{tikzpicture}[x=1mm, y=1mm]
  \def\R{40}
  \draw[bkGray, line width=0.6pt, -{Stealth[length=4pt]}] (-6,0) -- (46,0) node[anchor=west, font=\sffamily\footnotesize] {$x$};
  \draw[bkGray, line width=0.6pt, -{Stealth[length=4pt]}] (0,-4) -- (0,46) node[anchor=south, font=\sffamily\footnotesize] {$y$};
  \draw[bkInk!35, line width=0.8pt] (-3:\R) arc (-3:95:\R);
  % もとの道すじ（うすく）
  \draw[bkBlue!35, line width=1.4pt] (0,0) -- ({\R*cos(40)},0);
  \draw[bkOrange!40, line width=1.4pt] ({\R*cos(40)},0) -- (40:\R);
  \fill[bkInk!40] (40:\R) circle (0.9);
  \node[font=\sffamily\footnotesize, text=bkInk!60, anchor=west] at ($(40:\R)+(0.8,-0.6)$) {Q（角 $b$）};
  % 回した道すじ
  \draw[bkBlue, line width=1.8pt] (0,0) -- (30:{\R*cos(40)});
  \draw[bkOrange, line width=1.8pt] (30:{\R*cos(40)}) -- (70:\R);
  \draw[bkInk, line width=0.5pt] ($(30:{\R*cos(40)})+(210:2.6)$) -- ++(120:2.6) -- ++(30:2.6);
  \draw[bkInk, line width=0.8pt] (0,0) -- (70:\R);
  \fill[bkInk] (70:\R) circle (1.1);
  \node[font=\sffamily\footnotesize, anchor=south west, inner sep=1pt] at ($(70:\R)+(0.6,0.4)$) {R（角 $a+b$）};
  \node[font=\sffamily\footnotesize, text=bkBlue!80!black, anchor=north west, inner sep=1pt] at ($(30:15)+(0.8,-0.6)$) {長さ $\cos b$};
  \node[font=\sffamily\footnotesize, text=bkOrange!80!black, anchor=west, inner sep=1pt] at (22.6,28.6) {長さ $\sin b$};
  \draw[bkInk, line width=0.5pt] (8,0) arc (0:30:8);
  \node[font=\sffamily\footnotesize] at (15:10.5) {$a$};
  \draw[bkInk, line width=0.5pt] (30:12) arc (30:70:12);
  \node[font=\sffamily\footnotesize] at (55:14.5) {$b$};
  % 右の説明
  \node[font=\sffamily\footnotesize, anchor=north west, align=left, text width=76mm] at (50,44) {%
  \textcolor{bkBlue!80!black}{青い道}：長さ $\cos b$、向きは角 $a$\\
  \quad 横に $\cos b\times\cos a$\\
  \quad たてに $\cos b\times\sin a$\\[4pt]
  \textcolor{bkOrange!80!black}{オレンジの道}：長さ $\sin b$、\\
  \quad 向きは、青い道を90°回した角 $a+90°$\\
  \quad 横に $\sin b\times\cos(a+90°)=-\sin a\sin b$\\
  \quad たてに $\sin b\times\sin(a+90°)=\cos a\sin b$};
\end{tikzpicture}
\bkfigcaption{Q への道すじ（うすい色）を、まるごと角 $a$ だけ回すと、R への道すじ（こい色）になる。図は $a=30°$、$b=40°$。}
\end{bkfigure}

長さ $r$ で角 $\theta$ の向きに進むと、横に $r\cos\theta$、たてに $r\sin\theta$ 進みます（4-8②）。オレンジの道で使った2つの式は、4-8⑤で作った90°回す式です。
\bkivjline{$\cos(a+90°)=-\sin a$\qquad $\sin(a+90°)=\cos a$}

Rの座標は、2本の道の横どうし、たてどうしをたしたものです。
\begin{bkcenter}
\begin{bktable}{colspec={Q[l,m,wd=17mm] Q[c,m,wd=27mm] Q[c,m,wd=27mm] Q[c,m,wd=46mm]}, rowsep=2pt, colsep=3pt, row{1}={font=\sffamily\small}}
 & 青い道 & オレンジの道 & Rの座標（合計） \\
横（$x$） & $\cos a\cos b$ & $-\sin a\sin b$ & $\cos a\cos b-\sin a\sin b$ \\
たて（$y$） & $\sin a\cos b$ & $\cos a\sin b$ & $\sin a\cos b+\cos a\sin b$ \\
\end{bktable}
\end{bkcenter}
Rの座標は $(\cos(a+b),\ \sin(a+b))$ でもあったので、
\bkivjframe[110mm]{\large $\sin(a+b)=\sin a\cos b+\cos a\sin b$\\[3pt] $\cos(a+b)=\cos a\cos b-\sin a\sin b$}
これを\textbf{加法定理}といいます。\bkidx{かほうていり}{加法定理}

\bkivjhead{数で確かめる}
$a=30°$、$b=60°$ では、$\sin90°=1$ のはずです。
\begin{align*}
\sin30°\cos60°+\cos30°\sin60°&=0.5\times0.5+0.866\times0.866\\
&=0.25+0.75=1
\end{align*}
$\cos90°=0$ のはずです。
\bkivjline{$\cos30°\cos60°-\sin30°\sin60°=0.866\times0.5-0.5\times0.866=0$}
どちらも合っています。

\bkivjhead{ひき算のとき}
$b$ のかわりに $-b$ を入れます。4-8③の $\cos(-b)=\cos b$、$\sin(-b)=-\sin b$ を使うと、
\begin{align*}
\sin(a-b)&=\sin a\cos(-b)+\cos a\sin(-b)=\sin a\cos b-\cos a\sin b\\
\cos(a-b)&=\cos a\cos(-b)-\sin a\sin(-b)=\cos a\cos b+\sin a\sin b
\end{align*}
4つの式をまとめると、次のようになります。$\pm$（プラス・マイナス）を使うと、たし算とひき算の式を1つに書けます。$\sin(a\pm b)=\sin a\cos b\pm\cos a\sin b$ のように書き、上の符号どうし、下の符号どうしを組にして読みます。$\mp$ は、$\pm$ と上下が反対の記号で、$\cos(a\pm b)=\cos a\cos b\mp\sin a\sin b$ と書きます。
\begin{bkcenter}
\begin{bktable}{colspec={Q[c,m,wd=62mm] Q[c,m,wd=62mm]}, rowsep=3pt, colsep=3pt, cells={font=\small}}
たし算（$a+b$） & ひき算（$a-b$） \\
$\sin(a+b)=\sin a\cos b+\cos a\sin b$ & $\sin(a-b)=\sin a\cos b-\cos a\sin b$ \\
$\cos(a+b)=\cos a\cos b-\sin a\sin b$ & $\cos(a-b)=\cos a\cos b+\sin a\sin b$ \\
\end{bktable}
\bktabcaption{$\sin$ は「$\sin\cos$ と $\cos\sin$」で、符号は左と同じ。$\cos$ は「$\cos\cos$ と $\sin\sin$」で、符号は左と反対。}
\end{bkcenter}

\bksubsection{3}{もう1つの作り方：距離を2通りに計算する}

\leavevmode\label{bk:4-9:kyori}加法定理は、これからの巻で何度も使う大事な式なので、別の方法でも確かめておきます。単位円の上に、角 $a$ の点P $(\cos a,\ \sin a)$ と、角 $b$ の点Q $(\cos b,\ \sin b)$ をとります（$a$ のほうが大きいとします）。PとQの距離を、2通りの方法で計算します。

\noindent\begin{minipage}[c]{0.38\linewidth}
\centering
\begin{tikzpicture}[x=0.85mm, y=0.85mm]
  \def\R{20}
  \bkivjaxes
  \draw[bkInk!45, line width=0.8pt] (0,0) circle (\R);
  \fill[bkTheme!12] (0,0) -- (110:\R) -- (25:\R) -- cycle;
  \draw[bkInk, line width=0.8pt] (0,0) -- (110:\R) (0,0) -- (25:\R);
  \draw[bkRed, line width=1.5pt] (110:\R) -- (25:\R);
  \fill[bkInk] (110:\R) circle (1) (25:\R) circle (1);
  \node[font=\sffamily\footnotesize, anchor=south east, inner sep=1pt] at ($(110:\R)+(-0.4,0.4)$) {P};
  \node[font=\sffamily\footnotesize, anchor=south west, inner sep=1pt] at ($(25:\R)+(0.4,0.4)$) {Q};
  \draw[bkInk, line width=0.5pt] (25:6) arc (25:110:6);
  \node[font=\sffamily\footnotesize, fill=bkTheme!12, inner sep=0.5pt] at (68:9.6) {$a-b$};
  \node[font=\sffamily\footnotesize, anchor=north east, inner sep=1pt] at ($(110:11)+(-0.6,0)$) {1};
  \node[font=\sffamily\footnotesize, anchor=north west, inner sep=1pt] at ($(25:11)+(0.4,-0.6)$) {1};
\end{tikzpicture}
\end{minipage}\hfill
\begin{minipage}[c]{0.58\linewidth}
\raggedright
\textbf{方法1：座標で（4-4の2点の距離）}\par
\vspace{-2pt}
\[ \text{PQ}^2=(\cos a-\cos b)^2+(\sin a-\sin b)^2 \]
\textbf{方法2：三角形OPQで（4-7の余弦定理）}\par
OPとOQは半径なので1。そのあいだの角は $a-b$。
\[ \text{PQ}^2=1^2+1^2-2\times1\times1\times\cos(a-b) \]
\end{minipage}\par\vspace{6pt}
\noindent 方法1の式を、3-8の展開 $(p-q)^2=p^2-2pq+q^2$ で1つずつ広げます。2行めから3行めでは、4-8の $\sin^2\theta+\cos^2\theta=1$ を使います。
\begin{align*}
\text{PQ}^2&=(\cos^2a-2\cos a\cos b+\cos^2b)+(\sin^2a-2\sin a\sin b+\sin^2b)\\
&=(\cos^2a+\sin^2a)+(\cos^2b+\sin^2b)-2(\cos a\cos b+\sin a\sin b)\\
&=1+1-2(\cos a\cos b+\sin a\sin b)\\
&=2-2(\cos a\cos b+\sin a\sin b)
\end{align*}
方法2の式は $2-2\cos(a-b)$ です。同じ長さの2乗なので、2つは等しくなります。両辺から2をひき、$-2$ でわると、
\bkivjline{$\cos(a-b)=\cos a\cos b+\sin a\sin b$}
②で作った $\cos(a-b)$ の式と、ぴったり同じです。ちがう道で考えても、同じ式にたどりつきました。

\begin{seikaku}
$a-b$ が180°より大きいときは、角POQは $360°-(a-b)$ になります。でも、$\cos(360°-x)$ と $\cos x$ は同じ値なので、式は同じです。P、O、Qが一直線にならんで三角形ができないときも、方法1の計算で、同じ式が成り立つことが確かめられます。だから、加法定理は、どんな角 $a$、$b$ でも成り立ちます。
\end{seikaku}

\bksubsection{4}{2倍角の式}

\leavevmode\label{bk:4-9:nibai}加法定理で、$a$ も $b$ も同じ角 $\theta$ にしてみます。$a+b=2\theta$ になります。
\begin{align*}
\sin2\theta&=\sin\theta\cos\theta+\cos\theta\sin\theta=2\sin\theta\cos\theta\\
\cos2\theta&=\cos\theta\cos\theta-\sin\theta\sin\theta=\cos^2\theta-\sin^2\theta
\end{align*}
$\cos2\theta$ の式は、4-8の $\sin^2\theta+\cos^2\theta=1$ を使って、$\sin$ だけ、または $\cos$ だけの形にも書けます。
\begin{itemize}[leftmargin=1.2\zw, itemsep=2pt, topsep=2pt, label={\color{\bkTT}●}]
\item $\cos^2\theta=1-\sin^2\theta$ を入れると、$\cos2\theta=(1-\sin^2\theta)-\sin^2\theta=1-2\sin^2\theta$
\item $\sin^2\theta=1-\cos^2\theta$ を入れると、$\cos2\theta=\cos^2\theta-(1-\cos^2\theta)=2\cos^2\theta-1$
\end{itemize}
\bkivjframe[110mm]{$\sin2\theta=2\sin\theta\cos\theta$\\[3pt] $\cos2\theta=\cos^2\theta-\sin^2\theta$\\[2pt] $=1-2\sin^2\theta=2\cos^2\theta-1$}
これを\textbf{2倍角の式}といいます。\bkidx{にばいかく}{2倍角}$\theta=30°$ で確かめます。
\begin{align*}
\sin60°&=0.866 & 2\sin30°\cos30°&=2\times0.5\times0.866=0.866\\
\cos60°&=0.5 & 1-2\sin^230°&=1-2\times0.25=0.5
\end{align*}
どちらも合っています。

$\sin2\theta$ は、$2\sin\theta$ ではありません。$\theta=30°$ なら、$\sin60°=0.866$ ですが、$2\sin30°=1$ です。4-8⑦のグラフで見たように、$\sin2\theta$ は「周期が半分」、$2\sin\theta$ は「たてに2倍」で、まったくちがう波です。

\bkivjhead{2乗を、2倍の角に直す}
$\cos2\theta=1-2\sin^2\theta$ を、$\sin^2\theta$ について解きます（2-14）。$2\sin^2\theta$ を左へ、$\cos2\theta$ を右へ移すと、$2\sin^2\theta=1-\cos2\theta$。両辺を2でわって、
\bkivjframe[110mm]{$\sin^2\theta=\dfrac{1-\cos2\theta}{2}$\qquad $\cos^2\theta=\dfrac{1+\cos2\theta}{2}$}
（$\cos^2\theta$ の式は、$\cos2\theta=2\cos^2\theta-1$ を同じように解いて作ります。）

この式は、「2乗」という計算を、「2倍の角の $\cos$」に直します。2乗のままでは扱いにくい式も、この形にすると、4-8で学んだ波のグラフの動かし方で考えられます。

\noindent\begin{minipage}[c]{0.52\linewidth}
\begin{tikzpicture}
\begin{axis}[bkplot, width=68mm, height=42mm, xmin=0, xmax=395, ymin=-1.15, ymax=1.3, xtick={90,180,270,360}, xticklabels={}, ytick={-1,0.5,1}, yticklabels={$-1$,$\frac12$,1}, xlabel={$\theta$}, tick label style={font=\sffamily\footnotesize}]
  \addplot[bkthin, domain=0:360, samples=120] {sin(x)};
  \addplot[bkcurve, domain=0:360, samples=120] {sin(x)^2};
  \addplot[bkOrange, line width=0.8pt, dash pattern=on 2.5pt off 1.5pt, domain=0:360] {0.5};
  \node[font=\sffamily\footnotesize, anchor=north, inner sep=1pt] at (axis cs:90,-1.17) {90°};
  \node[font=\sffamily\footnotesize, anchor=north, inner sep=1pt] at (axis cs:180,-1.17) {180°};
  \node[font=\sffamily\footnotesize, anchor=north, inner sep=1pt] at (axis cs:270,-1.17) {270°};
  \node[font=\sffamily\footnotesize, anchor=north, inner sep=1pt] at (axis cs:360,-1.17) {360°};
\end{axis}
\end{tikzpicture}
\end{minipage}\hfill
\begin{minipage}[c]{0.45\linewidth}
\raggedright
青い線が $y=\sin^2\theta$、灰色の線が $y=\sin\theta$ です。$\sin^2\theta=\dfrac12-\dfrac12\cos2\theta$ なので、$\sin^2\theta$ のグラフは、周期180°の波が、高さ $\dfrac12$ の線（オレンジ）を中心にゆれている形です。2乗するとマイナスがなくなるので、0と1のあいだでゆれます。
\end{minipage}



\bksubsection{5}{噴水は、なぜ45°でいちばん遠くへ届く？}

\leavevmode\label{bk:4-9:funsui}3-9で、「同じ速さでふき出す水は、ななめ45°に向けたとき、いちばん遠くまで届く」と学びました。そのわけを、2倍角の式で説明します（空気の抵抗は考えません）。

\noindent\begin{minipage}[c]{0.56\linewidth}
\raggedright
水のつぶが、速さ $v$（m/s）で、地面から角 $\theta$ の向きにふき出します。4-6で斜面の重力を2つの向きに分けたように、ななめの速さも、横とたての2つに分けて考えます。長さ $v$ で角 $\theta$ の向きの矢印は、横に $v\cos\theta$、たてに $v\sin\theta$ の大きさをもちます（4-8の $(r\cos\theta,\ r\sin\theta)$ と同じ。くわしくは4-11）。
\end{minipage}\hfill
\begin{minipage}[c]{0.4\linewidth}
\centering
\begin{tikzpicture}[x=1mm, y=1mm]
  \draw[bkGray, line width=0.6pt] (-2,0) -- (40,0);
  \draw[bkBlue, line width=1.6pt, -{Stealth[length=5pt]}] (0,0) -- (35:30);
  \draw[bkOrange, line width=1.4pt, -{Stealth[length=4pt]}] (0,0) -- ({30*cos(35)},0);
  \draw[bkOrange, line width=1.4pt, -{Stealth[length=4pt]}] ({30*cos(35)},0) -- ({30*cos(35)},{30*sin(35)});
  \node[font=\sffamily\footnotesize, text=bkBlue!80!black, anchor=south east, inner sep=1pt] at (35:16) {$v$};
  \node[font=\sffamily\footnotesize, text=bkOrange!80!black, anchor=north, inner sep=1pt] at ({15*cos(35)},-0.8) {横 $v\cos\theta$};
  \node[font=\sffamily\footnotesize, text=bkOrange!80!black, anchor=west, inner sep=1pt] at ({30*cos(35)+0.8},8.6) {たて $v\sin\theta$};
  \draw[bkInk, line width=0.5pt] (7,0) arc (0:35:7);
  \node[font=\sffamily\footnotesize] at (17:10) {$\theta$};
\end{tikzpicture}
\end{minipage}

\begin{enumerate}[label={\bkcircled{\arabic*}}, leftmargin=2.4\zw, itemsep=3pt, topsep=4pt]
\item \textbf{たての動き}\quad 3-9の投げ上げと同じで、$t$ 秒後の高さは $y=(v\sin\theta)\,t-5t^2$（m）。5は、落ちる距離 $5t^2$ の5です（2-13・3-9。地球の重力の加速度\qty{9.8}{m/s^2}の半分を、まるめた数）。
\item \textbf{地面にもどる時間}\quad $y=0$ になる時間を求めます。$(v\sin\theta)\,t=5t^2$ の両辺を $t$ でわって（0でない $t$ を考える。2-14）、$v\sin\theta=5t$。だから $t=\dfrac{v\sin\theta}{5}$。
\item \textbf{横に進む距離}\quad 横には、同じ速さ $v\cos\theta$ で進むので、$x=v\cos\theta\times t$。これに②の $t$ を入れて、$x=\dfrac{v^2\sin\theta\cos\theta}{5}$。
\item \textbf{2倍角で書きかえる}\quad $\sin2\theta=2\sin\theta\cos\theta$ から、$\sin\theta\cos\theta=\dfrac{\sin2\theta}{2}$。これを入れると、
\end{enumerate}
\bkivjframe[70mm]{\large $x=\dfrac{v^2\sin2\theta}{10}$}
$v$ が同じなら、$x$ は $\sin2\theta$ に比例します。$\sin2\theta$ がいちばん大きい1になるのは、$2\theta=90°$、つまり $\theta=45°$ のときです。これが、45°でいちばん遠くへ届くわけです。分母の10は、①の5の2倍です（重力の加速度を $g$ と書くと、$x=\dfrac{v^2\sin2\theta}{g}$）。

\begin{bkcenter}
\begin{bktable}{colspec={Q[l,m,wd=34mm] *{5}{Q[c,m,wd=14mm]}}, rowsep=1.5pt}
角 $\theta$ & 15° & 30° & 45° & 60° & 75° \\
$\sin2\theta$ & 0.5 & 0.866 & 1 & 0.866 & 0.5 \\
届く距離（m） & 5 & 8.66 & 10 & 8.66 & 5 \\
\end{bktable}
\bktabcaption{$v=\qty{10}{m/s}$ のときの届く距離 $x=\dfrac{100\sin2\theta}{10}=10\sin2\theta$。}
\end{bkcenter}

\noindent\begin{minipage}[c]{0.5\linewidth}
\centering
\begin{tikzpicture}[x=1mm, y=1mm]
  % 速さ 10 m/s、角度 30°・45°・60°。1 m を 5.5 mm でかいた（落ちる距離 5t²）
  \fill[bkBlue!15] (-3,-2.5) rectangle (60,0);
  \draw[bkBlue!70!black, line width=0.6pt] (-3,0) -- (60,0);
  \foreach \ang/\col in {30/bkGreen!70!black, 45/bkBlue, 60/bkRed} {
    \pgfmathsetmacro\Rx{10*sin(2*\ang)}
    \draw[\col, line width=1.3pt, domain=0:\Rx, samples=40] plot ({5.5*\x},{5.5*(\x*tan(\ang)-5*\x*\x/(100*cos(\ang)*cos(\ang)))});}
  \fill[bkGray!80] (-1.5,0) rectangle (1.5,1.5);
  \node[font=\sffamily\footnotesize, text=bkRed!80!black, anchor=south, inner sep=1pt] at (5.5*4.33,5.5*3.75+0.6) {60°};
  \node[font=\sffamily\footnotesize, text=bkBlue!80!black, anchor=south, inner sep=1pt] at (5.5*5,5.5*2.5+0.6) {45°};
  \node[font=\sffamily\footnotesize, text=bkGreen!60!black, anchor=south, inner sep=1pt] at (5.5*5.8,5.5*1.15+0.4) {30°};
\end{tikzpicture}
\end{minipage}\hfill
\begin{minipage}[c]{0.46\linewidth}
\raggedright
30°と60°の届く距離は同じです。$\sin60°$ と $\sin120°$ は等しいからです（4-8③の $\sin(180°-\theta)=\sin\theta$）。15°と75°も同じです。実際の球技では、空気の抵抗や、投げる高さがあるので、いちばん遠くへ飛ぶ角は45°より少し小さくなります。
\end{minipage}

\bksubsection{6}{2つの波を1つにまとめる：合成}

\leavevmode\label{bk:4-9:gousei}やってみようの $\sin\theta+\sin(\theta+60°)$ にもどります。加法定理で $\sin(\theta+60°)$ を広げると、
\bkivjline{$\sin(\theta+60°)=\sin\theta\cos60°+\cos\theta\sin60°=0.5\sin\theta+0.866\cos\theta$}
だから、
\bkivjline{$\sin\theta+\sin(\theta+60°)=1.5\sin\theta+0.866\cos\theta$}
このように、$a\sin\theta+b\cos\theta$（$a$、$b$ は数）の形になります。この形の波を、1つの $\sin$ の波 $R\sin(\theta+\alpha)$ に書き直す方法を考えます。$\alpha$（アルファ）は、波のずれを表す角です。

\bkivjhead{加法定理を逆向きに使う}
$R\sin(\theta+\alpha)$ を加法定理で広げると、
\bkivjline{$R\sin(\theta+\alpha)=R\cos\alpha\sin\theta+R\sin\alpha\cos\theta$}
これが $a\sin\theta+b\cos\theta$ と同じになるには、$\sin\theta$ の前の数どうし、$\cos\theta$ の前の数どうしが等しければよいので、
\bkivjline{$R\cos\alpha=a$\qquad $R\sin\alpha=b$}

\noindent\begin{minipage}[c]{0.56\linewidth}
\raggedright
これは、「点 $(a,\ b)$ は、原点から距離 $R$、角 $\alpha$ の所にある」ということです（4-8②：半径 $r$ の円の上の点は $(r\cos\theta,\ r\sin\theta)$）。だから、
\begin{itemize}[leftmargin=1.2\zw, itemsep=2pt, topsep=2pt, label={\color{\bkTT}●}]
\item $R$ は、点 $(a,\ b)$ と原点の距離：$R=\sqrt{a^2+b^2}$（4-4）
\item $\alpha$ は、点 $(a,\ b)$ の動径の角
\end{itemize}
となります。
\end{minipage}\hfill
\begin{minipage}[c]{0.4\linewidth}
\centering
\begin{tikzpicture}[x=1mm, y=1mm]
  \draw[bkGray, line width=0.6pt, -{Stealth[length=4pt]}] (-3,0) -- (40,0) node[anchor=west, font=\sffamily\footnotesize] {$x$};
  \draw[bkGray, line width=0.6pt, -{Stealth[length=4pt]}] (0,-3) -- (0,24) node[anchor=south, font=\sffamily\footnotesize] {$y$};
  \draw[bkBlue, line width=1.5pt] (0,0) -- (30,0);
  \draw[bkOrange, line width=1.5pt] (30,0) -- (30,17.32);
  \draw[bkInk, line width=1pt] (0,0) -- (30,17.32);
  \fill[bkInk] (30,17.32) circle (1);
  \node[font=\sffamily\footnotesize, anchor=south west, inner sep=1pt] at (30.6,17.6) {$(a,\ b)$};
  \node[font=\sffamily\footnotesize, text=bkBlue!80!black, anchor=north, inner sep=1pt] at (15,-0.8) {$a$};
  \node[font=\sffamily\footnotesize, text=bkOrange!80!black, anchor=west, inner sep=1pt] at (30.8,8.7) {$b$};
  \node[font=\sffamily\footnotesize, anchor=south east, inner sep=1pt] at (14.4,9.2) {$R$};
  \draw[bkInk, line width=0.5pt] (8,0) arc (0:30:8);
  \node[font=\sffamily\footnotesize] at (15:10.6) {$\alpha$};
\end{tikzpicture}
\end{minipage}

\bkivjframe[110mm]{$a\sin\theta+b\cos\theta=\sqrt{a^2+b^2}\,\sin(\theta+\alpha)$\\[2pt]{\small $\alpha$ は、点 $(a,\ b)$ の動径の角（$\cos\alpha=\dfrac{a}{\sqrt{a^2+b^2}}$、$\sin\alpha=\dfrac{b}{\sqrt{a^2+b^2}}$）}}
このように、$a\sin\theta+b\cos\theta$ を1つの $\sin$ にまとめることを、三角関数の\textbf{合成}（ごうせい）といいます。\bkidx{ごうせい}{合成（三角関数）}

\bkivjhead{やってみようの答え}
$a=1.5$、$b=0.866$（$\frac{\sqrt3}{2}$）です。
\bkivjline{$R=\sqrt{1.5^2+0.866^2}=\sqrt{2.25+0.75}=\sqrt3\approx1.73$}$\cos\alpha=\dfrac{1.5}{1.73}=0.866$ なので、$\alpha=30°$ です（$\sin\alpha=\dfrac{0.866}{1.73}=0.5$ も合う）。だから
\bkivjline{$\sin\theta+\sin(\theta+60°)=\sqrt3\,\sin(\theta+30°)$}
たした波は、振幅 $\sqrt3\approx1.73$ の波で、$\sin\theta$ より30°早く進んでいます（4-8⑦：$\theta+30°$ は左へ30°）。山は $\theta+30°=90°$ の $\theta=60°$。表とグラフで見た「山の高さ1.73、山の位置60°」と、ぴったり合いました。

\bkivjhead{$b$ がマイナスのとき}
$\sin\theta-\cos\theta$ なら、$a=1$、$b=-1$ です。点 $(1,\ -1)$ は、原点から $\sqrt2$、角 $-45°$ の所にあるので、$\sin\theta-\cos\theta=\sqrt2\sin(\theta-45°)$ です。$\theta=45°$ で確かめると、左は $0.707-0.707=0$、右は $\sqrt2\sin0°=0$。合っています。

同じ周期の波は、いくつたしても、同じ周期の1つの波になります。たした波の振幅は、ふつう、振幅の和より小さくなります。山の位置がずれているからです。電気の回路では、ずれた電圧や電流をたす計算に、合成を毎日のように使います（6-10）。

\bksubsection{7}{積と和の式：うなりのしくみ}

\leavevmode\label{bk:4-9:unari}こんどは、周期が少しちがう2つの波をたします。そのために、加法定理から、もう1組の式を作ります。

\bkivjhead{積を和に直す式}
$\sin(a+b)$ と $\sin(a-b)$ の加法定理をたすと、$\cos a\sin b$ の部分が打ち消しあいます。
\begin{align*}
\sin(a+b)&=\sin a\cos b+\cos a\sin b\\
\sin(a-b)&=\sin a\cos b-\cos a\sin b\\
\sin(a+b)+\sin(a-b)&=2\sin a\cos b
\end{align*}
同じように、$\cos(a-b)$ から $\cos(a+b)$ をひくと、$\cos a\cos b$ の部分が打ち消しあって、
\bkivjline{$\cos(a-b)-\cos(a+b)=2\sin a\sin b$}
となります。左は和（や差）、右は積（かけ算）です。

\bkivjhead{和を積に直す式}
$\sin(a+b)+\sin(a-b)=2\sin a\cos b$ を、逆向きに読みます。$a+b=A$、$a-b=B$ とおくと、2つの式をたして $2a=A+B$、ひいて $2b=A-B$ なので、$a=\dfrac{A+B}{2}$、$b=\dfrac{A-B}{2}$。これを入れると、
\bkivjframe[100mm]{$\sin A+\sin B=2\sin\dfrac{A+B}{2}\cos\dfrac{A-B}{2}$}
2つの波の和が、「2つの角のまん中の $\sin$」と「2つの角の差の半分の $\cos$」の積になりました。

\bkivjhead{2つの音を重ねる}
音は、空気の圧力の波です。440\,Hz の音は、1秒に440回ゆれる波で、$t$ 秒後の大きさは $\sin(360°\times440t)$ と書けます（4-8⑨：1秒に440回転する点の高さ。1回転は360°）。440\,Hz と442\,Hz の音を重ねます。和を積に直す式で、$A=360°\times442t$、$B=360°\times440t$ とすると、
\begin{align*}
\frac{A+B}{2}&=360°\times\frac{442+440}{2}\,t=360°\times441t\\
\frac{A-B}{2}&=360°\times\frac{442-440}{2}\,t=360°\times1\times t
\end{align*}
なので、
\bkivjline{$\sin A+\sin B=2\cos(360°\times1\times t)\times\sin(360°\times441t)$}
\begin{itemize}[leftmargin=1.2\zw, itemsep=2pt, topsep=2pt, label={\color{\bkTT}●}]
\item $\sin(360°\times441t)$ は、441\,Hz の速い波。わたしたちには、2つの音のまん中の高さの音に聞こえます。
\item その前の $2\cos(360°\times1\times t)$ は、1秒に1回の、とてもゆっくりした波。速い波の振幅（音の大きさ）を、ゆっくり変えます。
\end{itemize}

\begin{bkfigure}
\begin{tikzpicture}
\begin{axis}[bkplot, width=120mm, height=46mm, xmin=0, xmax=1.08, ymin=-2.4, ymax=2.8, xtick={0.25,0.5,0.75,1}, xticklabels={}, ytick={-2,2}, xlabel={$t$（s）}, tick label style={font=\sffamily\footnotesize}]
  \addplot[bkBlue, line width=0.9pt, domain=0:1, samples=600] {sin(360*12*x)+sin(360*10*x)};
  \addplot[bkOrange, line width=1pt, dash pattern=on 2.5pt off 1.5pt, domain=0:1, samples=120] {2*cos(360*x)};
  \addplot[bkOrange, line width=1pt, dash pattern=on 2.5pt off 1.5pt, domain=0:1, samples=120] {-2*cos(360*x)};
  \node[font=\sffamily\footnotesize, anchor=north, inner sep=1pt] at (axis cs:0.25,-2.42) {0.25};
  \node[font=\sffamily\footnotesize, anchor=north, inner sep=1pt] at (axis cs:0.5,-2.42) {0.5};
  \node[font=\sffamily\footnotesize, anchor=north, inner sep=1pt] at (axis cs:0.75,-2.42) {0.75};
  \node[font=\sffamily\footnotesize, anchor=north, inner sep=1pt] at (axis cs:1,-2.42) {1};
  \node[font=\sffamily\footnotesize, anchor=south, inner sep=1pt] at (axis cs:0,2.1) {大};
  \node[font=\sffamily\footnotesize, anchor=south, inner sep=1pt] at (axis cs:0.5,2.1) {大};
  \node[font=\sffamily\footnotesize, anchor=south, inner sep=1pt] at (axis cs:1,2.1) {大};
  \node[font=\sffamily\footnotesize, anchor=south, inner sep=1pt] at (axis cs:0.25,2.1) {小};
  \node[font=\sffamily\footnotesize, anchor=south, inner sep=1pt] at (axis cs:0.75,2.1) {小};
\end{axis}
\end{tikzpicture}
\bkfigcaption{10\,Hz と12\,Hz の波の和（青）。見やすいように、ゆっくりした波でかいた。オレンジの点線は $\pm2\cos(360°\times t)$。音は、点線のはばが広い所で大きく、せまい所で小さくなる。}
\end{bkfigure}

\begin{yattemiyou}[表計算で、うなりを作ろう]
\begin{enumerate}[label={\bkcircled[bkGreen]{\arabic*}}, leftmargin=2.4\zw, itemsep=1pt, topsep=0pt]
\item A列に、時間 $t$ を0から1まで、0.01秒ずつ入れよう（101行）。
\item B列に \texttt{=SIN(2*PI()*10*A2)}、C列に \texttt{=SIN(2*PI()*12*A2)} と入れよう。10\,Hz と12\,Hz の波です（\texttt{2*PI()} は1回転の $2\pi$\,rad）。
\item D列に \texttt{=B2+C2} と入れて、2つの波をたそう。D列をグラフにして、大きい所と小さい所が、1秒に何回来るか数えよう。
\end{enumerate}
\begin{bkcenter}
\bksheet{4}{22}{%
1 & $t$（s） & 10\,Hz の波 & 12\,Hz の波 & たした値 \\
2 & 0 & 0 & 0 & 0 \\
3 & 0.01 & 0.588 & 0.685 & 1.272 \\
4 & 0.02 & 0.951 & 0.998 & 1.949 \\
26 & 0.24 & 0.588 & $-0.685$ & $-0.097$ \\
28 & 0.26 & $-0.588$ & 0.685 & 0.097 \\
53 & 0.51 & 0.588 & 0.685 & 1.272 \\}
\bktabcaption{0秒と0.5秒のあたりでは、2つの波の山がそろって、たした値は2近くになる。0.25秒のあたりでは山と谷が打ち消しあって、0近くになる。}
\end{bkcenter}
\end{yattemiyou}

\noindent 0秒、0.5秒、1秒のあたりで大きく、0.25秒、0.75秒のあたりで小さくなります。1秒に2回、つまり $12-10=2$ 回です。2つの波は、0秒では山がそろっていますが、12\,Hz の波が少しずつ先に進み、0.25秒では半周（山と谷）ずれて打ち消しあいます。0.5秒では、12\,Hz の波がちょうど1周多く進んで、また山がそろいます。このようすを式で表したのが、上の「和を積に直す式」です。

\bkivjhead{\bkhatenaref の答え}
音の大きさは、振幅 $2\cos(360°\times t)$ の「大きさ」（プラスかマイナスかは関係ない）で決まります。$\cos(360°\times t)$ は、$t=0$ で1、$t=0.5$ で $-1$、$t=1$ で1。大きさが最大の2になるのは、0秒、0.5秒、1秒、……と、\textbf{1秒に2回}です。これがウワンウワンの正体です。
\bkivjline{1秒あたりのうなりの回数＝$442-440=2$（回）}
一般に、振動数が $f_1$ と $f_2$ の音を重ねると、1秒に $f_1-f_2$ 回（大きいほうから小さいほうをひく）うなります。差の半分の振動数の $\cos$ は、大きさ（プラスかマイナスかは考えない）が、1秒にその2倍の回数だけ最大になるからです。2つの音の高さが近づくほど、うなりはゆっくりになり、同じ高さでは0回、つまりうなりは消えます。\bkidx{うなり}{うなり}

\begin{motto}[4]
\textbf{たくさんの波をたすと、どんな形も作れる}\par
周期のちがう波をたしていくと、なめらかでない形も作れます。$\sin\theta$ に、$\dfrac13\sin3\theta$、$\dfrac15\sin5\theta$、$\dfrac17\sin7\theta$、……とたしていくと、山の上が平らになり、四角い波（方形波）に近づきます。$\theta=90°$ で計算すると、$1-\dfrac13+\dfrac15-\dfrac17\approx0.724$ で、たす数をふやすと約0.785に近づきます。
\begin{bkcenter}
\begin{tikzpicture}
\begin{axis}[bkplot, width=118mm, height=34mm, xmin=0, xmax=395, ymin=-1.15, ymax=1.2, xtick={90,180,270,360}, xticklabels={}, ytick={-1,1}, xlabel={$\theta$}, tick label style={font=\sffamily\footnotesize}]
  \addplot[bkthin, domain=0:360, samples=120] {sin(x)};
  \addplot[bkcurve, domain=0:360, samples=400] {sin(x)+sin(3*x)/3+sin(5*x)/5+sin(7*x)/7};
  \addplot[bkOrange, line width=0.9pt, dash pattern=on 3pt off 1.5pt] coordinates {(0,0.785) (180,0.785) (180,-0.785) (360,-0.785)};
  \node[font=\sffamily\footnotesize, anchor=north, inner sep=1pt] at (axis cs:180,-1.17) {180°};
  \node[font=\sffamily\footnotesize, anchor=north, inner sep=1pt] at (axis cs:360,-1.17) {360°};
\end{axis}
\end{tikzpicture}
\end{bkcenter}
（灰色が $\sin\theta$、青が4つの波の和、オレンジの点線が方形波。）デジタル回路の0と1の信号も、たくさんの $\sin$ の和で表せます。この考え方を\textbf{フーリエ級数}といい、第9巻（9-1）で学びます。
\end{motto}

\bksubsection{8}{交流の電力と、$\sin^2$ の平均}

\leavevmode\label{bk:4-9:denryoku}家のコンセントの交流では、電圧は $V_0\sin\omega t$ の形で、プラスとマイナスを行ったり来たりします（4-8）。電気ストーブのような抵抗に流れる電流も、同じ形 $I_0\sin\omega t$ で行ったり来たりします（$V_0$、$I_0$ は山の高さ）。電力（1秒あたりに使う電気のエネルギー）は、電圧×電流なので、
\bkivjline{$p=V_0\sin\omega t\times I_0\sin\omega t=V_0I_0\sin^2\omega t$}
④の $\sin^2\theta=\dfrac12-\dfrac12\cos2\theta$ を使うと、
\bkivjline{$p=\dfrac{V_0I_0}{2}-\dfrac{V_0I_0}{2}\cos2\omega t$}

\noindent\begin{minipage}[c]{0.52\linewidth}
\begin{tikzpicture}
\begin{axis}[bkplot, width=68mm, height=46mm, xmin=0, xmax=395, ymin=-0.35, ymax=1.2, xtick={90,180,270,360}, xticklabels={}, ytick={0.5,1}, yticklabels={$\frac12$,1}, xlabel={$\omega t$}, tick label style={font=\sffamily\footnotesize}]
  \addplot[bkcurve, domain=0:360, samples=120] {sin(x)^2};
  \addplot[bkcurve2, line width=1.2pt, domain=0:360, samples=120] {sin(x)*sin(x-60)};
  \addplot[bkBlue, line width=0.7pt, dash pattern=on 2pt off 1.5pt, domain=0:360] {0.5};
  \addplot[bkOrange, line width=0.7pt, dash pattern=on 2pt off 1.5pt, domain=0:360] {0.25};
  \node[font=\sffamily\footnotesize, anchor=north, inner sep=1pt] at (axis cs:180,-0.37) {180°};
  \node[font=\sffamily\footnotesize, anchor=north, inner sep=1pt] at (axis cs:360,-0.37) {360°};
\end{axis}
\end{tikzpicture}
\end{minipage}\hfill
\begin{minipage}[c]{0.45\linewidth}
\raggedright
青い線は $\sin^2\omega t$（電力を $V_0I_0$ でわったもの）。高さ $\dfrac12$ の点線を中心に、周期180°でゆれています。山の $\dfrac12$ より上の部分を切りとって、谷にうめると、ちょうど高さ $\dfrac12$ の平らな形になります。オレンジの線は、あとで考える「電流が60°おくれる」ときの電力です。
\end{minipage}

\noindent 第2の部分 $-\dfrac{V_0I_0}{2}\cos2\omega t$ は、0を中心に、プラスとマイナスに同じだけゆれる波です。1周期で平均すると、プラスとマイナスが打ち消しあって0になります。だから、\textbf{電力の平均は $\dfrac{V_0I_0}{2}$}。$\sin^2$ の平均が $\dfrac12$ だということです。

数でも確かめます。$\sin^2$ の値を30°ごとに12こ（0°〜330°）求めると、0、0.25、0.75、1、0.75、0.25、0、0.25、0.75、1、0.75、0.25。合計は6なので、平均は $6\div12=0.5$ です（積分を使った確かめは5-15）。

\bkivjhead{実効値}
山の高さを $V_0=\sqrt2\,V_{\mathrm{e}}$、$I_0=\sqrt2\,I_{\mathrm{e}}$ と書きます。平均の電力 $\dfrac{V_0I_0}{2}$ に入れると、$\sqrt2\times\sqrt2=2$ なので、
\bkivjline{$\dfrac{\sqrt2V_{\mathrm{e}}\times\sqrt2I_{\mathrm{e}}}{2}=\dfrac{2V_{\mathrm{e}}I_{\mathrm{e}}}{2}=V_{\mathrm{e}}I_{\mathrm{e}}$}
となり、直流の「電圧×電流」と同じ形になります。この $V_{\mathrm{e}}$、$I_{\mathrm{e}}$ が\textbf{実効値}です。3-10で学んだ「\qty{100}{V}は実効値で、山の高さは約\qty{141}{V}」の $\sqrt2$ は、$\sin^2$ の平均の $\dfrac12$ から出てきたのです。

\bkivjhead{電流がおくれるとき}
モーターのようにコイル（巻いた電線）をふくむ機械では、電流の波が電圧の波よりおくれます。電流が60°おくれて $I_0\sin(\omega t-60°)$ になると、電力は次のようになります。
\bkivjline{$p=V_0I_0\sin\omega t\sin(\omega t-60°)$}⑦の「積を和に直す式」$2\sin a\sin b=\cos(a-b)-\cos(a+b)$ で、$a=\omega t$、$b=\omega t-60°$ とすると、$a-b=60°$、$a+b=2\omega t-60°$ なので、
\bkivjline{$p=\dfrac{V_0I_0}{2}\bigl\{\cos60°-\cos(2\omega t-60°)\bigr\}$}
$\cos(2\omega t-60°)$ の部分は、やはり平均すると0です。だから平均の電力は $\dfrac{V_0I_0}{2}\cos60°=\dfrac{V_0I_0}{2}\times0.5$ で、おくれがないときの半分になります（図のオレンジの線の平均 $\frac14$）。一般に、平均の電力は $V_{\mathrm{e}}I_{\mathrm{e}}\times\cos$（おくれの角）です。この $\cos$（おくれの角）を\textbf{力率}（りきりつ）といいます。\bkidx{りきりつ}{力率}電流が同じでも、力率が小さいと、使える電力は小さくなります。










\begin{taisetsu}
\def\tsl#1{\makebox[5.4\zw][l]{\sffamily\bfseries\small\color{\bkTT}#1}}
\tsl{加法定理}$\sin(a\pm b)=\sin a\cos b\pm\cos a\sin b$\\[3pt]
\tsl{}$\cos(a\pm b)=\cos a\cos b\mp\sin a\sin b$\\[3pt]
\tsl{2倍角}$\sin2\theta=2\sin\theta\cos\theta$、$\cos2\theta=1-2\sin^2\theta=2\cos^2\theta-1$\\[3pt]
\tsl{}$\sin^2\theta=\dfrac{1-\cos2\theta}{2}$、$\cos^2\theta=\dfrac{1+\cos2\theta}{2}$\\[3pt]
\tsl{合成}$a\sin\theta+b\cos\theta=\sqrt{a^2+b^2}\,\sin(\theta+\alpha)$\\[3pt]
\tsl{和を積に}$\sin A+\sin B=2\sin\dfrac{A+B}{2}\cos\dfrac{A-B}{2}$
\end{taisetsu}

\begin{senmonsho}[電気回路の教科書ふうの書き方]
\begin{bkhonbun}
正弦波交流電圧
\[ v=\sqrt2V\sin\omega t\ \text{〔V〕} \]
を負荷に加えたとき，電流
\[ i=\sqrt2I\sin(\omega t-\theta)\ \text{〔A〕} \]
が流れたとする．瞬時電力は
\[ p=vi=VI\cos\theta-VI\cos(2\omega t-\theta) \]
となる．第2項の1周期の平均は0であるから，平均電力は $P=VI\cos\theta$〔W〕である．$\cos\theta$ を力率という．
\end{bkhonbun}
\bkyaku{電圧の実効値を $V$、電流の実効値を $I$ と書いている（山の高さは $\sqrt2$ 倍）。$\theta$ は、電流が電圧よりおくれる角。電力 $p$ は、積を和に直す式で、一定の部分と、2倍の速さでゆれる部分に分かれる。ゆれる部分の平均は0なので、平均の電力は $VI\cos\theta$。}
\end{senmonsho}

\begin{kotoba}
\small
\bkword{加法定理}{2つの角の和や差の $\sin$・$\cos$ を、それぞれの角の $\sin$・$\cos$ で表す式。}
\bkword{2倍角の式}{加法定理で2つの角を同じにした式。$\sin2\theta$、$\cos2\theta$ を表す。}
\bkword{合成}{$a\sin\theta+b\cos\theta$ を、1つの $\sin$ の波 $R\sin(\theta+\alpha)$ にまとめること。}
\bkword{うなり}{振動数が少しちがう2つの音を重ねたとき、音の大きさがくり返し変わること。1秒に $f_1-f_2$ 回。}
\bkword{力率}{交流で、電圧と電流の波のおくれの角の $\cos$。平均の電力＝実効値の電圧×電流×力率。}
\end{kotoba}


\bkkigocards{\bkkigocard{$\alpha$}{アルファ}{合成した波のずれの角。点 $(a,\ b)$ の動径の角}{phase angle}{4-9}}
{\bkkigocard{$\pm$}{プラス・マイナス}{$+$ と $-$ の2つの式をまとめて書く記号}{plus-minus}{4-9}}






\begin{column}{モールの円：材料の中の力}
\noindent\begin{minipage}[c]{0.6\linewidth}
\raggedright
金属の棒を強さ $\sigma$（面積1あたりの力）で引っぱると、角 $\theta$ だけかたむいた切り口には、おしはなす力 $\sigma\cos^2\theta=\dfrac{\sigma}{2}+\dfrac{\sigma}{2}\cos2\theta$ と、ずらす力 $\sigma\sin\theta\cos\theta=\dfrac{\sigma}{2}\sin2\theta$ がはたらきます（2倍角の式）。この2つの組は、中心 $\bigl(\frac{\sigma}{2},\ 0\bigr)$、半径 $\frac{\sigma}{2}$ の円の上の点です（4-8②）。この円を\textbf{モールの円}といい、設計で使います。ずらす力は $\theta=45°$ で最大で、やわらかい金属の棒は45°の向きにすべります。
\end{minipage}\hfill
\begin{minipage}[c]{0.38\linewidth}
\centering
\begin{tikzpicture}[x=1mm, y=1mm]
  \draw[bkGray, line width=0.6pt, -{Stealth[length=4pt]}] (-4,0) -- (36,0);
  \node[anchor=north west, font=\sffamily\footnotesize, text=bkInk, align=left, inner sep=1pt] at (30.5,-1.2) {おしは\\なす力};
  \draw[bkGray, line width=0.6pt, -{Stealth[length=4pt]}] (0,-16) -- (0,18) node[anchor=south, font=\sffamily\footnotesize, text=bkInk] {ずらす力};
  \draw[bkBlue, line width=1.3pt] (14,0) circle (14);
  \fill[bkInk] (14,0) circle (0.8);
  \draw[bkInk, line width=0.8pt] (14,0) -- ++(60:14);
  \path (14,0) ++(60:14) coordinate (Q);
  \fill[bkOrange] (Q) circle (1.2);
  \draw[bkInk, line width=0.5pt] (20,0) arc (0:60:6);
  \node[font=\sffamily\footnotesize] at ($(14,0)+(30:9)$) {$2\theta$};
  \node[font=\sffamily\footnotesize, anchor=north, inner sep=1pt] at (14,-1.2) {$\frac{\sigma}{2}$};
  \node[font=\sffamily\footnotesize, anchor=south west, inner sep=1pt] at (28.6,0.8) {$\sigma$};
  \node[font=\sffamily\footnotesize, anchor=south west, inner sep=1pt] at ($(Q)+(0.8,0.6)$) {切り口の力};
\end{tikzpicture}
\end{minipage}
\end{column}

%% ---------------- 例題 ----------------
\bksection{reidai}{例題でたしかめよう}

\begin{reidai}[加法定理で値を求める]
次の値を求めましょう。\quad（1）$\sin75°$\quad（2）$\cos15°$\par
\kangaekata 三角定規の角（30°・45°・60°）の和や差に分けて、加法定理を使います。$\sqrt2\times\sqrt3=\sqrt6$（3-10）。\par
\shiki\hangindent=4.3\zw\hangafter=1 （1）$75°=45°+30°$\par
\bkivjb $\sin75°=\sin45°\cos30°+\cos45°\sin30°$\par
\bkivjb $\phantom{\sin75°}=\dfrac{\sqrt2}{2}\times\dfrac{\sqrt3}{2}+\dfrac{\sqrt2}{2}\times\dfrac12=\dfrac{\sqrt6}{4}+\dfrac{\sqrt2}{4}$\par
\bkivjq （2）$15°=45°-30°$\par
\bkivjb $\cos15°=\cos45°\cos30°+\sin45°\sin30°$\par
\bkivjb $\phantom{\cos15°}=\dfrac{\sqrt2}{2}\times\dfrac{\sqrt3}{2}+\dfrac{\sqrt2}{2}\times\dfrac12=\dfrac{\sqrt6}{4}+\dfrac{\sqrt2}{4}$\par
\bkivjb 確かめ：$\dfrac{2.449+1.414}{4}=\dfrac{3.863}{4}\approx0.966$\par
\bkivjb 電卓の $\sin75°=0.966$ と同じ。\par
\kotae\hangindent=4.3\zw\hangafter=1 （1）$\dfrac{\sqrt6+\sqrt2}{4}\approx0.966$\quad（2）$\dfrac{\sqrt6+\sqrt2}{4}\approx0.966$
\end{reidai}
$\sin75°$ と $\cos15°$ が同じになるのは、$75°+15°=90°$ だからです（4-6の $\sin(90°-\theta)=\cos\theta$）。

\begin{jibun}
次の値を求めましょう。\quad（1）$\cos75°$\quad（2）$\sin15°$
\bkjans{（1）\nobreak $\cos45°\cos30°-\sin45°\sin30°=\dfrac{\sqrt6-\sqrt2}{4}\approx0.259$。（2）\nobreak $\sin45°\cos30°-\cos45°\sin30°=\dfrac{\sqrt6-\sqrt2}{4}\approx0.259$。}
\end{jibun}

\begin{reidai}[$\sin a$、$\cos b$ から $\sin(a+b)$ を求める]
$a$、$b$ は0°と90°のあいだの角で、$\sin a=\dfrac35$、$\cos b=\dfrac{5}{13}$ です。$\sin(a+b)$ と $\cos(a+b)$ を求めましょう。\par
\kangaekata 加法定理には $\cos a$ と $\sin b$ も必要です。$\sin^2+\cos^2=1$（4-8）で求めます。0°〜90°の角なので、$\sin$ も $\cos$ も正。\par
\shiki\hangindent=2.6\zw\hangafter=1 $\cos^2a=1-\Bigl(\dfrac35\Bigr)^2=1-\dfrac{9}{25}=\dfrac{16}{25}$ より $\cos a=\dfrac45$\par
\bkivja $\sin^2b=1-\Bigl(\dfrac{5}{13}\Bigr)^2=1-\dfrac{25}{169}=\dfrac{144}{169}$ より $\sin b=\dfrac{12}{13}$\par
\bkivja $\sin(a+b)=\dfrac35\times\dfrac{5}{13}+\dfrac45\times\dfrac{12}{13}=\dfrac{15}{65}+\dfrac{48}{65}=\dfrac{63}{65}$\par
\bkivja $\cos(a+b)=\dfrac45\times\dfrac{5}{13}-\dfrac35\times\dfrac{12}{13}=\dfrac{20}{65}-\dfrac{36}{65}=-\dfrac{16}{65}$\par
\kotae\hangindent=2.6\zw\hangafter=1 $\sin(a+b)=\dfrac{63}{65}$、$\cos(a+b)=-\dfrac{16}{65}$
\end{reidai}
$\cos(a+b)$ が負なので、$a+b$ は90°より大きい角です。確かめに、$\Bigl(\dfrac{63}{65}\Bigr)^2+\Bigl(\dfrac{16}{65}\Bigr)^2$ を計算すると、$\dfrac{3969+256}{4225}=1$ になり、$\sin^2+\cos^2=1$ も成り立っています。

\begin{jibun}
$a$、$b$ は0°と90°のあいだの角で、$\sin a=\dfrac{5}{13}$、$\cos b=\dfrac45$ です。$\sin(a+b)$ と $\cos(a+b)$ を求めましょう。
\bkjans{$\cos a=\dfrac{12}{13}$、$\sin b=\dfrac35$。$\sin(a+b)=\dfrac{20}{65}+\dfrac{36}{65}=\dfrac{56}{65}$、$\cos(a+b)=\dfrac{48}{65}-\dfrac{15}{65}=\dfrac{33}{65}$。}
\end{jibun}

\begin{reidai}[加法定理で式を短くする]
$\sin(\theta+120°)+\sin(\theta-120°)$ を、加法定理で短くしましょう。\par
\kangaekata 2つを加法定理で広げてからたします。$\cos120°=-\dfrac12$、$\sin120°=\dfrac{\sqrt3}{2}$。\par
\shiki\hangindent=2.6\zw\hangafter=1 $\sin(\theta+120°)=\sin\theta\cos120°+\cos\theta\sin120°$\par
\bkivja $\phantom{\sin(\theta+120°)}=-\dfrac12\sin\theta+\dfrac{\sqrt3}{2}\cos\theta$\par
\bkivja $\sin(\theta-120°)=\sin\theta\cos120°-\cos\theta\sin120°$\par
\bkivja $\phantom{\sin(\theta-120°)}=-\dfrac12\sin\theta-\dfrac{\sqrt3}{2}\cos\theta$\par
\bkivja たすと、$\cos\theta$ の部分が打ち消しあって、\par
\bkivja $-\dfrac12\sin\theta-\dfrac12\sin\theta=-\sin\theta$\par
\bkivja 確かめ：$\theta=90°$ なら、$\sin210°+\sin(-30°)$ で、\par
\bkivja $-0.5-0.5=-1=-\sin90°$\par
\kotae\hangindent=2.6\zw\hangafter=1 $-\sin\theta$
\end{reidai}

\begin{jibun}
$\cos(\theta+120°)+\cos(\theta-120°)$ を、加法定理で短くしましょう。
\bkjans{$\cos(\theta+120°)=-\dfrac12\cos\theta-\dfrac{\sqrt3}{2}\sin\theta$、$\cos(\theta-120°)=-\dfrac12\cos\theta+\dfrac{\sqrt3}{2}\sin\theta$。たすと $-\cos\theta$。}
\end{jibun}

\noindent 例題3の答えを、グラフで確かめます。
\begin{bkfigure}
\begin{tikzpicture}
\begin{axis}[bkplot, width=120mm, height=40mm, xmin=0, xmax=395, ymin=-1.25, ymax=1.3, xtick={90,180,270,360}, xticklabels={}, ytick={-1,1}, xlabel={$\theta$}, tick label style={font=\sffamily\footnotesize}]
  \addplot[bkcurve2, line width=1.1pt, domain=0:360, samples=120] {sin(x+120)};
  \addplot[bkPurple, line width=1.1pt, domain=0:360, samples=120] {sin(x-120)};
  \addplot[bkcurve, domain=0:360, samples=120] {sin(x+120)+sin(x-120)};
  \node[font=\sffamily\footnotesize, anchor=north, inner sep=1pt] at (axis cs:90,-1.27) {90°};
  \node[font=\sffamily\footnotesize, anchor=north, inner sep=1pt] at (axis cs:180,-1.27) {180°};
  \node[font=\sffamily\footnotesize, anchor=north, inner sep=1pt] at (axis cs:270,-1.27) {270°};
  \node[font=\sffamily\footnotesize, anchor=north, inner sep=1pt] at (axis cs:360,-1.27) {360°};
\end{axis}
\end{tikzpicture}
\bkfigcaption{オレンジが $\sin(\theta+120°)$、紫が $\sin(\theta-120°)$、青が2つの和。青い波は、$\sin\theta$ を上下にひっくり返した $-\sin\theta$ になっている（90°で $-1$、270°で1）。}
\end{bkfigure}

\begin{reidai}[2倍角の式を使う]
（1）$\theta$ は0°と90°のあいだの角で、$\sin\theta=0.6$、$\cos\theta=0.8$ です。$\sin2\theta$ と $\cos2\theta$ を求めましょう。\par
（2）$\sin^215°$ を求めましょう（$\cos30°=0.866$）。\par
\kangaekata （1）は2倍角の式。（2）は $\sin^2\theta=\dfrac{1-\cos2\theta}{2}$ で、$2\theta=30°$ にします。\par
\shiki\hangindent=4.3\zw\hangafter=1 （1）$\sin2\theta=2\times0.6\times0.8=0.96$\par
\bkivjb $\cos2\theta=1-2\times0.6^2=1-2\times0.36=1-0.72=0.28$\par
\bkivjb 確かめ：$\cos^2\theta-\sin^2\theta=0.64-0.36=0.28$。同じ。\par
\bkivjq （2）$\sin^215°=\dfrac{1-\cos30°}{2}=\dfrac{1-0.866}{2}=\dfrac{0.134}{2}=0.067$\par
\kotae\hangindent=4.3\zw\hangafter=1 （1）$\sin2\theta=0.96$、$\cos2\theta=0.28$\quad（2）0.067
\end{reidai}

\begin{jibun}
（1）$\theta$ は0°と90°のあいだの角で、$\sin\theta=0.8$、$\cos\theta=0.6$ です。$\sin2\theta$ と $\cos2\theta$ を求めましょう。\par
（2）$\cos^215°$ を求めましょう（$\cos30°=0.866$）。
\bkjans{（1）\nobreak $\sin2\theta=2\times0.8\times0.6=0.96$、$\cos2\theta=1-2\times0.64=-0.28$。（2）\nobreak $\dfrac{1+0.866}{2}=0.933$。}
\end{jibun}

\noindent 2倍角の式 $\sin2\theta=2\sin\theta\cos\theta$ を、電卓でいろいろな角について確かめてみましょう（小数第4位まで）。
\begin{bkcenter}
\begin{bktable}{colspec={Q[c,m,wd=14mm] Q[c,m,wd=20mm] Q[c,m,wd=20mm] Q[c,m,wd=22mm] Q[c,m,wd=34mm]}, rowsep=1.5pt}
$\theta$ & $\sin\theta$ & $\cos\theta$ & $\sin2\theta$ & $2\sin\theta\cos\theta$ \\
10° & 0.1736 & 0.9848 & 0.3420 & 0.3420 \\
20° & 0.3420 & 0.9397 & 0.6428 & 0.6428 \\
40° & 0.6428 & 0.7660 & 0.9848 & 0.9848 \\
70° & 0.9397 & 0.3420 & 0.6428 & 0.6428 \\
\end{bktable}
\bktabcaption{どの角でも、右の2列は同じ値になる。}
\end{bkcenter}

\begin{reidai}[ななめに投げた物の届く距離]
水のつぶが、速さ\qty{10}{m/s}で、地面から30°の向きにふき出します。届く距離を求めましょう。また、同じ速さで、いちばん遠くへ届く角と、その距離を求めましょう（$\sin60°=0.866$）。\par
\kangaekata ⑤の $x=\dfrac{v^2\sin2\theta}{10}$ を使います。\par
\shiki\hangindent=2.6\zw\hangafter=1 $x=\dfrac{10^2\times\sin60°}{10}=\dfrac{100\times0.866}{10}=8.66$\par
\bkivja いちばん遠いのは $\sin2\theta=1$ のとき。$2\theta=90°$ で $\theta=45°$\par
\bkivja $x=\dfrac{100\times1}{10}=10$\par
\kotae\hangindent=2.6\zw\hangafter=1 約\qty{8.7}{m}。45°のとき\qty{10}{m}
\end{reidai}

\begin{jibun}
ボールを、速さ\qty{20}{m/s}で、地面から75°の向きに投げます（空気の抵抗は考えない）。届く距離と、いちばん遠くへ届くときの距離を求めましょう（$\sin150°=0.5$）。
\bkjans{$x=\dfrac{400\times0.5}{10}=\qty{20}{m}$。いちばん遠いのは45°のときで、$\dfrac{400}{10}=\qty{40}{m}$。}
\end{jibun}

\noindent\begin{minipage}[c]{0.5\linewidth}
\centering
\begin{tikzpicture}[x=1mm, y=1mm]
  % 速さ 20 m/s、角度 30°・45°・60°。1 m を 1.5 mm でかいた（落ちる距離 5t²）
  \fill[bkBlue!15] (-3,-2.5) rectangle (64,0);
  \draw[bkBlue!70!black, line width=0.6pt] (-3,0) -- (64,0);
  \foreach \ang/\col in {30/bkGreen!70!black, 45/bkBlue, 60/bkRed} {
    \pgfmathsetmacro\Rx{40*sin(2*\ang)}
    \draw[\col, line width=1.3pt, domain=0:\Rx, samples=50] plot ({1.5*\x},{1.5*(\x*tan(\ang)-5*\x*\x/(400*cos(\ang)*cos(\ang)))});}
  \node[font=\sffamily\footnotesize, text=bkRed!80!black, anchor=south, inner sep=1pt] at (26,23.1) {60°};
  \node[font=\sffamily\footnotesize, text=bkBlue!80!black, anchor=south, inner sep=1pt] at (31,15.6) {45°};
  \node[font=\sffamily\footnotesize, text=bkGreen!60!black, anchor=south, inner sep=1pt] at (24,8.1) {30°};
  \node[font=\sffamily\footnotesize, anchor=north east, inner sep=1pt] at (53,-3) {34.6 m};
  \node[font=\sffamily\footnotesize, anchor=north, inner sep=1pt] at (60.5,-3) {40 m};
  \draw[bkGray, line width=0.4pt] (52,-0.8) -- (52,0.8) (60,-0.8) -- (60,0.8);
\end{tikzpicture}
\end{minipage}\hfill
\begin{minipage}[c]{0.46\linewidth}
\raggedright
速さ\qty{20}{m/s}で、30°・45°・60°に投げたときの通り道です。30°と60°の届く距離は同じ約\qty{34.6}{m}ですが、空中にいる時間（⑤の②の $t=\dfrac{v\sin\theta}{5}$）は、30°では2.0秒、60°では約3.5秒です。高く上がるぶん、時間がかかります。
\end{minipage}

\begin{reidai}[三角関数の合成]
次の式を $R\sin(\theta+\alpha)$ の形にし、最大の値と、そのときの $\theta$（0°〜360°）を求めましょう。\par
（1）$\sin\theta+\sqrt3\cos\theta$\quad（2）$3\sin\theta+4\cos\theta$\par
\kangaekata 点 $(a,\ b)$ の、原点からの距離が $R$、動径の角が $\alpha$。最大は $\sin(\theta+\alpha)=1$ のとき、つまり $\theta+\alpha=90°$。\par
\shiki\hangindent=4.3\zw\hangafter=1 （1）点 $(1,\ \sqrt3)$。$R=\sqrt{1+3}=2$。$\cos\alpha=\dfrac12$、$\sin\alpha=\dfrac{\sqrt3}{2}$ より $\alpha=60°$\par
\bkivjb 最大は2。$\theta+60°=90°$ より $\theta=30°$\par
\bkivjq （2）点 $(3,\ 4)$。$R=\sqrt{9+16}=\sqrt{25}=5$。$\cos\alpha=\dfrac35$、$\sin\alpha=\dfrac45$\par
\bkivjb 電卓で $\cos^{-1}0.6=53.13\cdots$ より $\alpha\approx53.1°$\par
\bkivjb 最大は5。$\theta=90°-53.1°=36.9°$\par
\kotae\hangindent=4.3\zw\hangafter=1 （1）$2\sin(\theta+60°)$。最大2（$\theta=30°$）\par
\bkivjq （2）$5\sin(\theta+\alpha)$（$\alpha\approx53.1°$）。最大5（$\theta\approx36.9°$）
\end{reidai}
（2）の最大は、$3+4=7$ ではなく5です。$3\sin\theta$ の山（90°）と $4\cos\theta$ の山（0°）の位置がずれているので、山の高さはたし算になりません。

\begin{jibun}
次の式を $R\sin(\theta+\alpha)$ の形にし、最大の値と、そのときの $\theta$（0°〜360°）を求めましょう。\par
（1）$\sqrt3\sin\theta+\cos\theta$\quad（2）$4\sin\theta+3\cos\theta$
\bkjans{（1）\nobreak $R=2$、$\alpha=30°$。$2\sin(\theta+30°)$、最大2（$\theta=60°$）。（2）\nobreak $R=5$、$\cos\alpha=0.8$ より $\alpha\approx36.9°$。最大5（$\theta\approx53.1°$）。}
\end{jibun}

\begin{reidai}[うなりの回数]
（1）440\,Hz と443\,Hz の音を同時に鳴らすと、1秒に何回うなりますか。聞こえる音の高さは、およそ何Hzですか。\par
（2）440\,Hz のおんさと、ギターの弦を同時に鳴らすと、1秒に2回うなりました。弦の音は何Hzですか。\par
\kangaekata ⑦より、うなりの回数は振動数の差。聞こえる高さは、2つの振動数のまん中（和の半分）。\par
\shiki\hangindent=4.3\zw\hangafter=1 （1）$443-440=3$。$\dfrac{440+443}{2}=441.5$\par
\bkivjq （2）差が2なので、弦は440より2高いか、2低い。$440+2=442$、$440-2=438$\par
\kotae\hangindent=4.3\zw\hangafter=1 （1）1秒に3回、約441.5\,Hz\quad（2）442\,Hz か438\,Hz
\end{reidai}
（2）で、どちらかを決めるには、弦を少ししめて（音を高くして）みます。うなりがゆっくりになれば438\,Hz、速くなれば442\,Hz です。調律では、うなりが消えるまで弦を調節します。

\begin{jibun}
（1）262\,Hz と260\,Hz の音を同時に鳴らすと、1秒に何回うなりますか。聞こえる音の高さは、およそ何Hzですか。\par
（2）330\,Hz の音と、ある音を同時に鳴らすと、1秒に3回うなりました。ある音は何Hzですか。
\bkjans{（1）\nobreak $262-260=2$ で1秒に2回、$\dfrac{262+260}{2}=\qty{261}{Hz}$。（2）\nobreak 333\,Hz か327\,Hz。}
\end{jibun}

\begin{reidai}[交流の平均の電力]
実効値\qty{100}{V}の電圧で、実効値\qty{10}{A}の電流が流れます。平均の電力を求めましょう。\par
（1）電流が電圧とずれないとき\quad（2）電流が電圧より60°おくれるとき\par
\kangaekata ⑧より、平均の電力＝実効値の電圧×実効値の電流×$\cos$（おくれの角）。ずれないときは、おくれの角が0°で $\cos0°=1$。\par
\shiki\hangindent=4.3\zw\hangafter=1 （1）$100\times10\times\cos0°=100\times10\times1=1000$\par
\bkivjq （2）$100\times10\times\cos60°=1000\times0.5=500$\par
\kotae\hangindent=4.3\zw\hangafter=1 （1）\qty{1000}{W}\quad（2）\qty{500}{W}
\end{reidai}

\begin{jibun}
実効値\qty{100}{V}の電圧で、実効値\qty{5}{A}の電流が流れます。（1）電流が電圧とずれないとき、（2）電流が30°おくれるときの、平均の電力を求めましょう（$\cos30°=0.866$）。
\bkjans{（1）\nobreak $100\times5=\qty{500}{W}$。（2）\nobreak $500\times0.866=\qty{433}{W}$。}
\end{jibun}

%% ---------------- まちがえやすいところ・物理工学 ----------------
\bksection{tsumazuki}{まちがえやすいところをチェック}

\begin{tsumazuki}{かっこの中を、そのまま分けてしまう}
\bkng{$\sin(a+b)=\sin a+\sin b$}\bkok{$\sin(a+b)=\sin a\cos b+\cos a\sin b$}
$\sin$ は、かけ算ではないので、かっこの中のたし算を分けられません。$a=30°$、$b=60°$ で確かめると、左は1、まちがいの右は1.366です。
\end{tsumazuki}

\begin{tsumazuki}{$\cos(a+b)$ の符号}
\bkng{$\cos(a+b)=\cos a\cos b+\sin a\sin b$}\bkok{$\cos(a+b)=\cos a\cos b-\sin a\sin b$}
$\cos$ の加法定理は、かっこの中と符号が反対になります。$a=b=45°$ で確かめると、$\cos90°=0$。○の式は $0.5-0.5=0$ で合いますが、×の式は1になります。
\end{tsumazuki}

\begin{tsumazuki}{$\sin2\theta$ を $2\sin\theta$ にする}
\bkng{$\sin60°=2\sin30°=1$}\bkok{$\sin60°=2\sin30°\cos30°=0.866$}
角を2倍にしても、$\sin$ の値は2倍になりません。$\sin$ の値は1をこえないことからも、×はおかしいとわかります。
\end{tsumazuki}

\begin{tsumazuki}{合成の最大を、係数の和にする}
\bkng{$3\sin\theta+4\cos\theta$ の最大は $3+4=7$}\bkok{最大は $\sqrt{3^2+4^2}=5$}
2つの波の山は、ちがう角で来ます。山どうしは重ならないので、最大は和より小さくなります。
\end{tsumazuki}

\begin{butsuri}
\bkivjbox{楽器の調律}
\vspace{2pt}\noindent\begin{minipage}[c]{0.6\linewidth}
\raggedright
ピアノの調律師は、おんさや、となりの弦と鳴らし比べて、うなりを聞きながら弦の張りを調節します。うなりの回数は振動数の差なので、1秒に1回のうなりなら、差は1\,Hz。耳では区別しにくい小さなちがいも、うなりなら数えられます。
\end{minipage}\hfill
\begin{minipage}[c]{0.36\linewidth}
\centering
\begin{tikzpicture}[x=1mm, y=1mm]
  % おんさ
  \draw[bkGray!80!black, line width=1.6pt, line cap=round] (0,0) -- (0,8) (0,8) to[out=90, in=-90] (-3,12) -- (-3,26) (0,8) to[out=90, in=-90] (3,12) -- (3,26);
  \foreach \r in {4,7,10} {\draw[bkBlue, line width=0.7pt] (6+\r*0.3,{18-\r*0.6}) arc (-30:30:\r);}
  % ギター
  \filldraw[fill=bkOrange!35, draw=bkOrange!70!black, line width=0.6pt] (24,6) ellipse (7 and 6);
  \filldraw[fill=bkOrange!35, draw=bkOrange!70!black, line width=0.6pt] (24,15) ellipse (5 and 4.5);
  \fill[white] (24,9) circle (1.8);
  \draw[bkOrange!70!black, line width=0.6pt] (24,9) circle (1.8);
  \filldraw[fill=bkOrange!60!black, draw=bkOrange!70!black, line width=0.5pt] (23,19) rectangle (25,32);
  \draw[bkInk!70, line width=0.3pt] (23.5,2) -- (23.5,32) (24.5,2) -- (24.5,32);
\end{tikzpicture}
\end{minipage}
\bkivjbox{ノイズキャンセル}
ノイズキャンセル機能のあるイヤホンは、外の音をマイクで聞きとり、それと山と谷が反対の音を、スピーカーから出します。$\sin(\theta+180°)=-\sin\theta$（4-8）なので、2つをたすと $\sin\theta+(-\sin\theta)=0$ となり、雑音が消えます。電車や飛行機の、低くて一定の音ほどよく消えます。
\end{butsuri}

\begin{butsuri}
\bkivjbox{三相交流}
発電所から工場へ送る電気は、山の位置が120°ずつずれた3つの交流です（三相交流）。3つの電圧の和は、いつも0です。
\bkivjline{$\sin\theta+\sin(\theta-120°)+\sin(\theta-240°)=0$}
$\sin(\theta-240°)=\sin(\theta+120°)$（1周の360°をたした）なので、\bkreidairef{3}から、後ろの2つの和は $-\sin\theta$ です。だから全体は0。電線を6本でなく3本にでき、大きなモーターをなめらかに回せます。
\begin{bkcenter}
\begin{tikzpicture}
\begin{axis}[bkplot, width=112mm, height=28mm, xmin=0, xmax=395, ymin=-1.25, ymax=1.3, xtick={90,180,270,360}, xticklabels={}, ytick={-1,1}, xlabel={$\theta$}, tick label style={font=\sffamily\footnotesize}]
  \addplot[bkcurve, line width=1.2pt, domain=0:360, samples=120] {sin(x)};
  \addplot[bkcurve2, line width=1.2pt, domain=0:360, samples=120] {sin(x-120)};
  \addplot[bkPurple, line width=1.2pt, domain=0:360, samples=120] {sin(x-240)};
  \addplot[bkInk, line width=1.4pt, dash pattern=on 3pt off 1.5pt, domain=0:360] {0};
  \node[font=\sffamily\footnotesize, anchor=north, inner sep=1pt] at (axis cs:180,-1.27) {180°};
  \node[font=\sffamily\footnotesize, anchor=north, inner sep=1pt] at (axis cs:360,-1.27) {360°};
  \node[font=\sffamily\footnotesize, anchor=south, inner sep=1pt] at (axis cs:300,1.32) {3つの和は0（黒の点線）};
\end{axis}
\end{tikzpicture}
\end{bkcenter}
\end{butsuri}



%% ---------------- 章末問題 ----------------
\bkendproblems
\bkproblemgroup{基本}{4}

\bkproblem 加法定理を使って、次の値を求めましょう\mbox{（\bkreidairef{1}）}。\quad（1）$\sin105°$\quad（2）$\cos105°$
\bkpans{$105°=60°+45°$。（1）\nobreak $\dfrac{\sqrt3}{2}\cdot\dfrac{\sqrt2}{2}+\dfrac12\cdot\dfrac{\sqrt2}{2}=\dfrac{\sqrt6+\sqrt2}{4}\approx0.966$。（2）\nobreak $\dfrac12\cdot\dfrac{\sqrt2}{2}-\dfrac{\sqrt3}{2}\cdot\dfrac{\sqrt2}{2}=\dfrac{\sqrt2-\sqrt6}{4}\approx-0.259$。}

\bkproblem $a$、$b$ は0°と90°のあいだの角で、$\sin a=\dfrac{8}{17}$、$\cos b=\dfrac35$ です。$\sin(a+b)$ と $\cos(a+b)$ を求めましょう\mbox{（\bkreidairef{2}）}。
\bkpans{$\cos a=\dfrac{15}{17}$、$\sin b=\dfrac45$。$\sin(a+b)=\dfrac{24}{85}+\dfrac{60}{85}=\dfrac{84}{85}$、$\cos(a+b)=\dfrac{45}{85}-\dfrac{32}{85}=\dfrac{13}{85}$。}

\bkproblem （1）$\theta$ は0°と90°のあいだの角で、$\sin\theta=0.28$、$\cos\theta=0.96$ です。$\sin2\theta$ と $\cos2\theta$ を求めましょう\mbox{（\bkreidairef{4}）}。\\（2）$\sin^222.5°$ を求めましょう（$\cos45°=0.707$）。
\bkpans{（1）\nobreak $\sin2\theta=0.5376$、$\cos2\theta=1-2\times0.0784=0.8432$。（2）\nobreak $\dfrac{1-0.707}{2}\approx0.146$。}

\bkproblem $\sin\theta+\cos\theta$ を $R\sin(\theta+\alpha)$ の形にし、最大の値と、そのときの $\theta$（0°〜360°）を求めましょう\mbox{（\bkreidairef{6}）}。
\bkpans{点 $(1,\ 1)$ で $R=\sqrt2$、$\alpha=45°$。$\sqrt2\sin(\theta+45°)$。最大 $\sqrt2$（$\theta=45°$）。}

\bkproblem 速さ\qty{15}{m/s}で、地面から30°の向きにボールを投げます。届く距離と、いちばん遠くへ届くときの距離を求めましょう\mbox{（\bkreidairef{5}）}（$x=\dfrac{v^2\sin2\theta}{10}$、$\sin60°=0.866$）。
\bkpans{$\dfrac{225\times0.866}{10}\approx\qty{19.5}{m}$。45°のとき $\dfrac{225}{10}=\qty{22.5}{m}$。}

\bkproblem $\sin(\theta+60°)+\sin(\theta-60°)$ を、加法定理で短くしましょう\mbox{（\bkreidairef{3}）}。
\bkpans{$2\sin\theta\cos60°=\sin\theta$。}

\bkproblem 330\,Hz と333\,Hz の音を同時に鳴らすと、1秒に何回うなりますか。聞こえる音の高さは、およそ何Hzですか\mbox{（\bkreidairef{7}）}。
\bkpans{1秒に3回。約331.5\,Hz。}

\bkproblem $a$、$b$ は0°と90°のあいだの角で、$\sin a=\dfrac45$、$\cos b=\dfrac{12}{13}$ です。$\sin(a+b)$ と $\cos(a+b)$ を求めましょう\mbox{（\bkreidairef{2}）}。
\bkpans{$\cos a=\dfrac35$、$\sin b=\dfrac{5}{13}$。$\sin(a+b)=\dfrac{48}{65}+\dfrac{15}{65}=\dfrac{63}{65}$、$\cos(a+b)=\dfrac{36}{65}-\dfrac{20}{65}=\dfrac{16}{65}$。}

\bkproblem $\theta$ は90°と180°のあいだの角で、$\sin\theta=0.6$、$\cos\theta=-0.8$ です。$\sin2\theta$ と $\cos2\theta$ を求めましょう\mbox{（\bkreidairef{4}）}。
\bkpans{$\sin2\theta=2\times0.6\times(-0.8)=-0.96$、$\cos2\theta=1-2\times0.36=0.28$。}

\bkproblemgroup{チャレンジ}{4}

\bkproblem $5\sin\theta+12\cos\theta$ を $R\sin(\theta+\alpha)$ の形にし、最大の値と、そのときの $\theta$ を求めましょう\mbox{（\bkreidairef{6}）}（$\alpha$ は電卓で、小数第1位まで）。
\bkpans{$R=\sqrt{25+144}=13$。$\cos\alpha=\dfrac{5}{13}$ より $\alpha\approx67.4°$。最大13（$\theta\approx22.6°$）。}

\bkproblem $\sin\theta+\sin(\theta-120°)+\sin(\theta-240°)=0$ を、加法定理で確かめましょう\mbox{（\bkreidairef{3}）}（$\cos240°=-\dfrac12$、$\sin240°=-\dfrac{\sqrt3}{2}$）。
\bkpans{$\sin(\theta-120°)=-\dfrac12\sin\theta-\dfrac{\sqrt3}{2}\cos\theta$、$\sin(\theta-240°)=-\dfrac12\sin\theta+\dfrac{\sqrt3}{2}\cos\theta$。3つをたすと $\sin\theta-\sin\theta+0=0$。}

\bkproblem 実効値\qty{200}{V}の電圧で、実効値\qty{10}{A}の電流が流れます。電流が電圧より60°おくれるときの平均の電力と、おくれがないときの平均の電力を求めましょう\mbox{（\bkreidairef{8}）}。
\bkpans{$200\times10\times0.5=\qty{1000}{W}$。おくれがないとき\qty{2000}{W}。}

\bkproblem 440\,Hz のおんさとピアノの弦を同時に鳴らすと、1秒に4回うなりました。弦を少ししめて音を高くすると、うなりはゆっくりになりました。弦の音は何Hzでしたか\mbox{（\bkreidairef{7}）}。
\bkpans{差が4なので436\,Hz か444\,Hz。音を高くしてうなりがゆっくりになった（440\,Hz に近づいた）ので、436\,Hz。}

\begin{furikaeri}
\begin{checklist}
\item 加法定理を単位円で導き、$\sin(a\pm b)$・$\cos(a\pm b)$ を計算できる。\bkfurref{bk:4-9:kahou}
\item 2倍角の式と $\sin^2\theta=\dfrac{1-\cos2\theta}{2}$ を作り、使える。\bkfurref{bk:4-9:nibai}
\item 合成と「積と和の式」で、うなりや電力を説明できる。\bkfurref{bk:4-9:gousei}
\end{checklist}
\end{furikaeri}

\bkfukidashi{delta}{波は、たしても波。次は、向きと大きさをもつ矢印、ベクトルの世界へ出発だ！}
